
% ════════════════════════════════════════════
% 研赛摘要（00.摘要.tex）—— 标题“摘要”和标签“关键词”使用随模板携带的隶书字体。
% ════════════════════════════════════════════
\begingroup
\renewcommand*{\mriteabstracttitlefont}{\mriteli}
\renewcommand*{\mritekeywordslabelfont}{\mriteli}

\begin{abstract}

本文针对算力约束下提升大语言模型能力的资源配置建模问题，结合综合评价、标度律建模、约束优化与技术演进预测等方法，提出了从“数据质量与配比刻画”到“算力最优配置”再到“能力前沿预测”的系统解决方案。

\textbf{针对问题一：} 本问是综合评价与回归建模问题，需对 22 个质量指标预处理并建立综合评价模型、定义并消解质量冲突、建立 17 域配比与交叉熵损失的定量关系。首先，逐指标判定方向并统一为“越高越好”（其中广告字段经原文抽查确认为“第 1 类=无广告”的正向指标），对 8 个多维列表型指标用 softmax 期望压缩为标量，再作极差归一化、用熵权法客观赋权得到样本级、语料级与领域级质量评分；其次，以成对标准化得分差定义冲突、以 Kendall 协同系数检验一致性，并给出“多数指标方向裁决”的消解规则；再次，用岭正则线性混合回归刻画 17 域配比到 13 域损失的线性关系，并检验追加质量项的增量解释力；最后在 1M、60M、1B 三套检验集上验证并做跨尺度外推。结果表明，全量 $272{,}505$ 条质量记录（另审计出 19 处字段缺失并按中位数填充）的 Kendall 协同系数为 $0.069$、冲突率 $22.3\%$（低于独立基准 $28.9\%$），说明指标间为正但不强的一致关系；抽样集 A1 与扩展集 A2/A3 的域级质量高度吻合（arxiv $0.703$ 对 $0.701$，github $0.455$ 对 $0.455$），领域质量排序为 arxiv $>$ stackexchange $>$ c4 $>$ commoncrawl $>$ book $\approx$ wikipedia $>$ github；配比—损失模型训练平均 $R^2=0.629$、同尺度检验 $0.619$，60M 与 1B 跨尺度 Pearson 相关达 $0.782$ 与 $0.651$；追加混合质量指数后平均 $R^2$ 仅提升 $9.9\times10^{-5}$，说明配比本身已编码质量信息，故将质量作为独立资源交由后续问题处理。

\textbf{针对问题二：} 本问是标度律建模问题，需拟合经典标度律、将数据质量纳入得到广义标度律、分析边际效用与弹性并推导质量与参数规模的替代条件。首先，用稳健非线性最小二乘在 Pythia 主拟合数据上拟合经典形式并做模型留出验证；其次，在控制 $N$、$D$ 的组内条件下审计质量字段方向；再次，比较四类含质量项的广义标度律形式，以拟合优度与内插/外推表现择优；最后计算弹性、等损失曲线与替代率。结果表明，经典标度律为 $L=1.6898+0.3540N^{-0.3400}+1.2403D^{-0.2799}$，$R^2=0.9999998$；方向审计的组内相关为 $-0.925$，确认 $Q_{\text{score}}$ 是质量指标（$Q=1$ 最优）；最优形式为加性质量缺口 $L=E+AN^{-\alpha}+BD^{-\beta}+C(1-Q)^{\gamma}$，联合拟合 $R^2=0.9896$，参数 $E=1.6558$、$A=0.3808$、$\alpha=0.3266$、$B=1.2515$、$\beta=0.2767$、$C=0.3544$、$\gamma=0.9779$；B3 插值轨迹的 $D$ 幂指数为 $-0.2805$，与拟合值 $-0.2799$ 高度一致；参考点 $N=1$B、$D=100$B、$Q=0.6$ 处质量弹性为 $-0.084$，强于规模弹性 $-0.049$，质量提升 $0.1$ 等价于参数量增加约 $0.28$B（在 $N=12$B 处等价于约 $7.68$B），验证了“以质量换规模”的可行性。

\textbf{针对问题三：} 本问是约束优化问题，需在算力预算下联合优化参数量、数据量与数据质量，并分析预算跨越量级时的结构性转移。首先，严格按赛题公式构建 $C=6ND+D[g(Q)-g(Q_0)]_{+}+\eta N D L_{\text{ctx}}$ 的总成本模型，并解析得到临界上下文长度 $L_{\text{crit}}=6/\eta=30{,}000$ tokens；其次，以问题一的质量基线 $Q_0=0.551$ 与附件 C7 的可行上下文长度为基础，用 SLSQP 在对数尺度上联合优化并比较附录 B 的三种质量成本函数；最后扫描 $10^{18}$–$10^{25}$ FLOPs 预算，以“质量投入启动/饱和边界”与“质量提升完成度弹性”定量识别结构性转移。结果表明，三档预算 $10^{19}$、$10^{22}$、$10^{24}$ FLOPs 下最优配置为 $(N^*,D^*,Q^*)=(0.225\text{B},6.9\text{B},0.551)$、$(6.089\text{B},229.5\text{B},1.000)$、$(44.937\text{B},3418.0\text{B},1.000)$；预算路径被两个门槛切成三段制度——$C\lesssim1.78\times10^{19}$ 时完全不投资质量（份额 $0$）、$C\approx7.08\times10^{19}$ 时质量投入份额达峰值 $0.407$、$C\gtrsim1\times10^{20}$ 后质量饱和、份额单调回落至 $0.6\%$，构成明确的结构性转移；当 $L_{\text{ctx}}=32{,}768>L_{\text{crit}}$ 时注意力开销达训练开销的 $1.09$ 倍，最优参数量由 $6.089$B 压缩至 $4.318$B，最优损失上升约 $6.0\%$。质量成本函数形式只在最低预算档显著影响最优质量（$0.551\sim1.000$）。

\textbf{针对问题四：} 本问是技术演进建模与前沿预测问题，需结合历史评测数据分离规模扩张与非规模技术进步的贡献并预测未来 12 至 24 个月能力前沿。首先明确时间轴（提交日期）、开源口径（开源权重且许可允许研究复现）与模型类型口径（pretrained 与 chat/finetuned 分组），得到 $2{,}346$ 条开源筛选样本；其次完成附件 C8 的逐任务聚合（$1{,}860$ 个模型目录、$37$ 个子任务维度）并与 C1 交叉校验；再次按可比性分层建立 Loss—Benchmark 桥接，并用面板回归分离规模效应与年份效应；最后用 logistic 饱和模型与结构外推两条路线给出预测区间，并引用附件 C4 的算力与开源权重证据设定情景。结果表明，桥接全样本 $r=-0.510$（$R^2=0.26$），映射误差须显式讨论；规模系数为 $13.84$ 分/10 倍参数、非规模进步 $4.84$ 分/年；长周期（2019—2025）前沿提升中规模扩张贡献 $50.5\%$、非规模技术进步 $49.5\%$，而在评测标准一致的近端窗口内前沿提升的 $95.4\%$ 来自非规模因素；未来 12/24 个月开源模型能力前沿预测为 $56.3$（$90\%$ 置信区间 $[53.6,64.5]$）与 $56.6$（$[53.8,66.5]$），算力增长放缓情景下为 $53.3$ 与 $55.8$；逐任务上指令遵循已近饱和（剩余空间 $10.0\%$），而研究生级问答与多轮对话仍有 $72.3\%$ 与 $61.3\%$ 的剩余空间。

综上，本文完成了数据质量评价、广义标度律建模、算力约束最优配置与能力前沿预测的全链条研究，各问结论均给出可复现的数值证据与不确定性刻画，可为算力受限条件下大模型训练的资源分配、数据治理与技术路线选择提供理论依据和决策参考。

\keywords{大语言模型；标度律；数据质量；资源配置优化；算力约束；能力前沿预测}

\end{abstract}
\endgroup
